% ══════════════════════════════════════════════════════════════════════
% Chapter 2: The Signalling Problem
% ══════════════════════════════════════════════════════════════════════

\chapter{The Signalling Problem}
\label{ch:signalling}

\epigraph{``Words are also actions, and actions are a kind of words.''}{Ralph Waldo Emerson}

% ── Introduction ──────────────────────────────────────────────────────

In every human society, actors choose signals strategically: a politician
picks rhetoric, a suitor offers gifts, a priest designs liturgy.  The
sender's fundamental problem is that they cannot control how receivers
will \emph{interpret} the signal---each receiver composes the signal with
their own repertoire.

This chapter formalises the core game-theoretic structure that governs
strategic communication.  Following Spence~\cite{spence1973} and
Crawford \& Sobel~\cite{crawford1982}, we model the signalling problem
as a game in which a sender chooses a signal, a receiver interprets it
through their repertoire, and both receive payoffs from the resulting
interpretations.  The crucial departure from classical game theory is
that payoffs depend not on the signal itself but on what receivers
\emph{make of it}---how it composes with their existing repertoire of
ideas.

\medskip
\noindent\textbf{Chapter overview.}  We define the \emph{signalling game}
(\S\ref{sec:ch2:game}), establish the \emph{babbling equilibrium}
(\S\ref{sec:ch2:babbling}), analyse \emph{audience structure}
(\S\ref{sec:ch2:audience}), prove \emph{depth bounds} on outcomes
(\S\ref{sec:ch2:depth}), study \emph{resonance in game outcomes}
(\S\ref{sec:ch2:resonance}), and explore \emph{strategic properties}
(\S\ref{sec:ch2:strategic}).  All eighteen theorems are machine-verified.

% ── §1. The Signalling Game ──────────────────────────────────────────
\section{The Signalling Game}
\label{sec:ch2:game}

\begin{definition}[Signalling Game]
\label{def:signalling-game}
A \emph{signalling game} $G = (\Rep, \Pay_S, \Pay_R)$ consists of:
\begin{itemize}
  \item A receiver repertoire $\Rep \in \IS^*$;
  \item A sender payoff function $\Pay_S : \IS^* \to \mathbb{N}$;
  \item A receiver payoff function $\Pay_R : \IS^* \to \mathbb{N}$.
\end{itemize}
The payoff depends on the \emph{interpretation} (the list of composed
ideas), not the signal directly.
\end{definition}

\begin{lstlisting}[caption={SignallingGame in Lean~4},label=lst:signalling-game]
structure SignallingGame (I : Type*) [IdeaticSpace I] where
  repertoire : Repertoire I
  sender_payoff : List I -> Nat
  receiver_payoff : List I -> Nat
\end{lstlisting}

\begin{definition}[Outcome]
\label{def:outcome}
The \emph{outcome} of sending signal~$s$ in game~$G$ is
\[
  \operatorname{outcome}(G, s) = \operatorname{interpret}(G.\Rep, s).
\]
\end{definition}

\begin{definition}[Optimal Signal]
\label{def:optimal-signal}
A signal~$s$ is \emph{optimal} for the sender if no alternative yields
strictly higher payoff:
\[
  \forall s'.\; \Pay_S(\operatorname{outcome}(G, s'))
  \leq \Pay_S(\operatorname{outcome}(G, s)).
\]
\end{definition}

\begin{definition}[Constant Sender Payoff]
\label{def:constant-payoff}
A game has \emph{constant sender payoff} if
$\Pay_S(\mathbf{x}) = \Pay_S(\mathbf{y})$ for all interpretation
lists~$\mathbf{x}, \mathbf{y}$.
\end{definition}

% ── §2. Babbling Equilibrium ─────────────────────────────────────────
\section{Babbling Equilibrium}
\label{sec:ch2:babbling}

\begin{theorem}[Babbling Outcome]
\label{thm:babbling-outcome}
Sending void produces the original repertoire unchanged:
\[
  \operatorname{outcome}(G, \void) = G.\Rep.
\]
\end{theorem}

\begin{proof}[Proof sketch]
Unfold the definitions and apply Theorem~\ref{thm:silence-transparent}
from Chapter~1.  See \texttt{babbling\_outcome}.
\end{proof}

\begin{remark}
The ``babbling equilibrium'' of Crawford \& Sobel~\cite{crawford1982}
exists in every signalling game because silence is transparent---the
receiver's mind is left exactly as it was.
\end{remark}

\begin{theorem}[Constant Payoff Babbling Optimality]
\label{thm:constant-payoff}
If sender payoff is constant, every signal is optimal---including void.
\end{theorem}

\begin{proof}[Proof sketch]
If $\Pay_S$ is constant, $\Pay_S(\operatorname{outcome}(G,s'))
= \Pay_S(\operatorname{outcome}(G,s))$ for all $s,s'$, so the
optimality condition holds trivially.  See \texttt{constant\_payoff\_optimal}.
\end{proof}

\begin{remark}
When the sender genuinely does not care what the receiver thinks,
information transmission collapses entirely.  This is the cheapest of
cheap talk.
\end{remark}

\begin{theorem}[Outcome Length Invariance]
\label{thm:outcome-length}
Every signal produces the same number of interpretations:
\[
  |\operatorname{outcome}(G, s)| = |G.\Rep|.
\]
\end{theorem}

\begin{proof}[Proof sketch]
By Theorem~\ref{thm:conserves-count}.  See \texttt{outcome\_length}.
\end{proof}

\begin{remark}
You cannot make someone think \emph{more} thoughts by choosing a
cleverer signal.  The number of schemas activated is a fixed property of
the receiver, not the sender.
\end{remark}

\begin{theorem}[Singleton Repertoire Outcome]
\label{thm:singleton-outcome}
A receiver with a single idea~$r$ produces $[r \compose s]$.
\end{theorem}

\begin{proof}[Proof sketch]
Direct computation.  See \texttt{singleton\_outcome}.
\end{proof}

\begin{remark}
The fanatic (single-idea mind) has a perfectly predictable interpretation
of any signal.  Totalitarian receivers are, paradoxically, the easiest
to communicate with.
\end{remark}

% ── §3. Audience Structure ───────────────────────────────────────────
\section{Audience Structure}
\label{sec:ch2:audience}

\begin{theorem}[Merged Audience]
\label{thm:merged-audience}
Signalling to a combined audience equals the concatenation of separate
outcomes:
\[
  \operatorname{outcome}(\Rep_1 \mathbin{+\!+} \Rep_2, s)
  = \operatorname{outcome}(\Rep_1, s) \mathbin{+\!+}
    \operatorname{outcome}(\Rep_2, s).
\]
\end{theorem}

\begin{proof}[Proof sketch]
By \texttt{List.map\_append}.  See \texttt{merged\_audience\_outcome}.
\end{proof}

\begin{remark}
A political speech to a diverse audience (workers $\cup$ elites) produces
interpretations that are the disjoint union of what each subgroup would
have produced separately.  There is no ``synergistic'' interpretation
effect from mere co-presence.
\end{remark}

\begin{theorem}[Empty Audience]
\label{thm:empty-audience}
Signalling to no one produces nothing:
$\operatorname{outcome}([], s) = []$.
\end{theorem}

\begin{proof}[Proof sketch]
By definition.  See \texttt{empty\_audience}.
\end{proof}

\begin{remark}
If nobody is there to interpret, there are no interpretations.  A signal
without a receiver is not a communicative act---formalising Wittgenstein's
point that meaning requires use~\cite{wittgenstein1953}.
\end{remark}

\begin{theorem}[Blank Audience Transparency]
\label{thm:blank-audience}
A receiver whose repertoire is all void interprets any signal as
$n$~copies of the signal itself:
\[
  \operatorname{outcome}(\operatorname{replicate}(n, \void), s)
  = \operatorname{replicate}(n, s).
\]
\end{theorem}

\begin{proof}[Proof sketch]
By induction on~$n$, using $\void \compose s = s$ at each step.
See \texttt{blank\_audience}.
\end{proof}

% ── §4. Depth Bounds on Outcomes ─────────────────────────────────────
\section{Depth Bounds on Outcomes}
\label{sec:ch2:depth}

\begin{theorem}[Outcome Depth Bound]
\label{thm:outcome-depth-bound}
\[
  \operatorname{depthSum}(\operatorname{outcome}(G, s))
  \leq \operatorname{depthSum}(G.\Rep) + |G.\Rep| \cdot \depth(s).
\]
\end{theorem}

\begin{proof}[Proof sketch]
Direct application of the Total Hermeneutic Bound
(Theorem~\ref{thm:total-hermeneutic-bound}).
See \texttt{outcome\_depth\_bound}.
\end{proof}

\begin{remark}
The sender's capacity to inject complexity into the receiver's mind is
linearly bounded.  This is why propaganda targets breadth (many receivers)
rather than depth (one deep signal).
\end{remark}

\begin{theorem}[Void Signal Zero Cost]
\label{thm:void-zero-cost}
The void signal adds zero total depth to the outcome:
\[
  \operatorname{depthSum}(\operatorname{outcome}(G, \void))
  = \operatorname{depthSum}(G.\Rep).
\]
\end{theorem}

\begin{proof}[Proof sketch]
By Theorem~\ref{thm:silence-depthsum}.  See \texttt{void\_signal\_zero\_cost}.
\end{proof}

% ── §5. Resonance in Game Outcomes ───────────────────────────────────
\section{Resonance in Game Outcomes}
\label{sec:ch2:resonance}

\begin{theorem}[Resonant Repertoire Consensus]
\label{thm:resonant-repertoire}
If all elements in the repertoire resonate with a fixed idea~$a$, then
all interpretations of any signal also resonate with $a \compose s$.
\end{theorem}

\begin{proof}[Proof sketch]
For each $r \in \Rep$ with $r \res a$, by Theorem~\ref{thm:durkheimian-consensus}
we get $r \compose s \res a \compose s$.
See \texttt{resonant\_repertoire\_consensus}.
\end{proof}

\begin{remark}
A culturally homogeneous audience (everyone resonates with the same
archetype) produces interpretations that all resonate with that
archetype's interpretation.  Cultural consensus is self-reinforcing.
\end{remark}

\begin{theorem}[Resonant Signals Produce Resonant Outcomes]
\label{thm:resonant-signals-outcomes}
If $s_1 \res s_2$, then for each position~$i$, the interpretations
$r_i \compose s_1$ and $r_i \compose s_2$ resonate.
\end{theorem}

\begin{proof}[Proof sketch]
By induction on the repertoire, applying resonance-compose-left at each
step.  See \texttt{resonant\_signals\_resonant\_outcomes}.
\end{proof}

% ── §6. Strategic Properties ─────────────────────────────────────────
\section{Strategic Properties}
\label{sec:ch2:strategic}

\begin{theorem}[Self-Signalling]
\label{thm:self-signalling}
If $\mathit{sender\_idea} \in \Rep$, then
$\mathit{sender\_idea} \compose s \in \operatorname{outcome}(G, s)$.
\end{theorem}

\begin{proof}[Proof sketch]
The sender's idea is mapped to its composition with~$s$ by the
interpretation function.  See \texttt{self\_signalling}.
\end{proof}

\begin{remark}
An insider---someone whose perspective is represented in the audience---
can always produce resonant interpretations by sending signals resonant
with their own position.
\end{remark}

\begin{theorem}[Composition Signal]
\label{thm:composition-signal}
\[
  \operatorname{outcome}(\Rep, s_1 \compose s_2)
  = \operatorname{outcome}(\operatorname{interpret}(\Rep, s_1), s_2).
\]
\end{theorem}

\begin{proof}[Proof sketch]
By associativity of composition: for each~$r$,
$r \compose (s_1 \compose s_2) = (r \compose s_1) \compose s_2$.
See \texttt{composition\_signal}.
\end{proof}

\begin{remark}
A complex ritual (compound signal) has the same effect as first
``priming'' the audience with~$s_1$ and then delivering~$s_2$.  Ritual
structure exploits associativity.
\end{remark}

\begin{theorem}[Void Prepend]
\label{thm:void-prepend}
$\operatorname{outcome}(\Rep, \void \compose s)
 = \operatorname{outcome}(\Rep, s)$.
\end{theorem}

\begin{proof}[Proof sketch]
By void-right: $r \compose (\void \compose s) = r \compose s$ after
applying associativity and the void-right identity.
See \texttt{void\_prepend\_signal}.
\end{proof}

\begin{theorem}[Sequential Signal]
\label{thm:sequential-signal}
\[
  \operatorname{interpret}(\operatorname{interpret}(\Rep, s_1), s_2)
  = \operatorname{interpret}(\Rep, s_1 \compose s_2).
\]
\end{theorem}

\begin{proof}[Proof sketch]
By functoriality of map and associativity.
See \texttt{sequential\_signal}.
\end{proof}

\begin{remark}
Multi-stage rituals (initiation sequences, liturgical calendars) produce
the same interpretive effect as a single complex ritual.  Ceremonial
structure is algebraically decomposable.
\end{remark}

\begin{theorem}[Void Captures Signal]
\label{thm:void-captures}
If $\void \in G.\Rep$, then $s \in \operatorname{outcome}(G, s)$.
\end{theorem}

\begin{proof}[Proof sketch]
By Theorem~\ref{thm:negative-capability}.
See \texttt{void\_captures\_signal}.
\end{proof}

\begin{remark}
An ``open-minded'' receiver (one who retains $\void$ = pure receptivity)
always captures the sender's intended signal.
\end{remark}

\begin{theorem}[MaxDepth Outcome Bound]
\label{thm:max-outcome-depth}
\[
  \operatorname{maxInterpDepth}(\operatorname{outcome}(G, s))
  \leq \operatorname{maxInterpDepth}(G.\Rep) + \depth(s).
\]
\end{theorem}

\begin{proof}[Proof sketch]
By Theorem~\ref{thm:max-interp-depth}.
See \texttt{max\_outcome\_depth}.
\end{proof}

\begin{remark}
No matter how profound a signal, the deepest interpretation anyone can
produce is bounded by their existing depth plus the signal's depth.  The
receiver's capacity limits understanding, not the sender's intent---a
formalisation of ``you can't teach beyond the student's readiness.''
\end{remark}

\begin{theorem}[Babbling Payoff]
\label{thm:babbling-payoff}
\[
  \Pay_S(\operatorname{outcome}(G, \void))
  = \Pay_S(G.\Rep).
\]
\end{theorem}

\begin{proof}[Proof sketch]
By Theorem~\ref{thm:babbling-outcome}.
See \texttt{babbling\_payoff}.
\end{proof}

\begin{remark}
The babbling payoff equals the ``status quo'' payoff.  A signal is worth
sending only if it yields strictly more than silence.
\end{remark}

% ── Summary ──────────────────────────────────────────────────────────
\section*{Chapter Summary}

This chapter formalised the signalling problem within IDT\@.  The central
game-theoretic structure---a signalling game with repertoire-dependent
payoffs---was introduced, and the \emph{babbling equilibrium}
(Theorem~\ref{thm:babbling-outcome}) was established: silence always
leaves the receiver's mind unchanged.  We showed that outcomes have
invariant length (Theorem~\ref{thm:outcome-length}), bounded depth
(Theorem~\ref{thm:outcome-depth-bound}), and that resonant signals
produce resonant outcomes (Theorem~\ref{thm:resonant-signals-outcomes}).
The composition signal theorem (Theorem~\ref{thm:composition-signal})
reveals that compound signals decompose via associativity, providing the
algebraic foundation for analysing multi-stage rituals.

All eighteen theorems are verified in
\texttt{FormalAnthropology/IDT\_Signal\_Ch02.lean}.
